\documentclass[10pt]{article}

\def\Type{LessonCaelan}
\def\TypeNum{Lesson 0} %Lesson #
\def\TypeTitle{Modeling Change \& Uncertainty} %Lesson Title
\def\Semester{Fall 2021} %Not Used 
\def\Course{MATH 1190 \& 1210}%Course number and name
\def\Targets{ \bf\color{ProcessBlue}A1 }
\def\desmosClassCode{{\bf ZAJ$3$AZ}}

\usepackage{preambleDJ} 

%%%%%%%%%%%% BEGIN DOCUMENT %%%%%%%%%%%%%%%
%%%%%%%%%%%% BEGIN DOCUMENT %%%%%%%%%%%%%%%
%%%%%%%%%%%% BEGIN DOCUMENT %%%%%%%%%%%%%%%
%%%%%%%%%%%% BEGIN DOCUMENT %%%%%%%%%%%%%%%
%%%%%%%%%%%% BEGIN DOCUMENT %%%%%%%%%%%%%%%

\begin{document}
\pagestyle{otherpages}
%\thispagestyle{empty}
\thispagestyle{firstpage}


\vs{0.4}


\textbf{Intended Learning Outcomes}
After this lesson, you will be able to 

\begin{itemize}
    \item Meet your instructor and LA for the first time (\target{I0})
    \item Understand the key ideas, tools and methods used in the course. (\target{I0})
    \item Establish clear communication and learn about your support options. (\target{I0})
\end{itemize}

\vspace{-0.5cm}

\hrulefill

\,



\minil {\it Intro to the course}\\

\newpage 



{\bf\color{blue}Warm-up}\exercise \ltrm{\bf\color{Plum} I0}Pair up with a classmate and interview each other using the questions below. \\

    
\begin{enumerate}[label=\textbf{Q\arabic*.}]
    \item \textbf{What's your name?}
    
    \vspace{1.5cm} % Space for answer
    
    \item \textbf{What is your major? If you haven't chosen one yet, what are you considering?}
    
    \vspace{2.5cm}
    
    \item \textbf{What's on your playlist?}
    
    \vspace{2.5cm}
    
    \item \textbf{What was your favorite class in high school or at UVA?}
    
    \vspace{2.5cm}
    
    
    \item \textbf{What are your goals for this semester?}
    
    \vspace{2.5cm}
    
    \item \textbf{How do you approach studying new material? Share some learning techniques or tips that have worked for you in the past.}

    \vspace{2.5cm}

    \item \textbf{Which word would you use to describe mathematics? Calculus?}
    
\end{enumerate}





\newpage
\exercise{\bf\color{Plum} I0} \\
\mpageT{.4}{.6}[\hs{.1}]{
The Keeling Curve is the record of atmospheric CO2 in parts per million (ppm). This data set is named after Charles Keeling who started collecting this data at Mauna Loa Observatory, starting in 1958.  A sample of atmospheric CO$_2$ is shown below in five-year intervals from starting in 1965 through 2010.  Note that ``year" means ``years since 1965."  For example, $t=10$ represents the year 1975.
}{
\graphicsC{keelingDataMMA}{3.5}
}

\begin{center}
\begin{tabular}{|c|c|c|c|c|c|c|c|c|c|c|}
\hline
year & 0 & 5 & 10 & 15 & 20 & 25 & 30 & 35 & 40 & 45 \\
\hline
CO$_2$ level & 320 & 325 & 331 & 338 & 345 & 354 & 360 & 369 & 379 & 389 \\
\hline
\end{tabular}
\end{center}

{\bf How much  CO$_2$ is in the atmosphere in the year 2025?}

Let's start with some intuition.  What's your best guess for the amount of CO$_2$ in the atmosphere in the year 2025?  

\vspace{5cm}

On a scale of $0-100$, how confident are you in your estimate? 

\newpage

What decisions did you make and/or questions did you have when you made your best guess? 

\vfill

Work together to brainstorm at least two strategies to refine your estimate and/or increase your confidence in your estimate. 

\vfill

In what ways would you describe how the amount of CO$_2$ is changing each year?
\vfill

What other information might be helpful to further improve your estimate and/or build confidence in your estimate? What mathematical tools might be helpful?

\vfill

\end{document}